%=============================================================================
% ABSTRACT
% Flow (lecturer's feedback): background -> technical gap in existing work ->
% proposed framework -> key results -> code link.
%=============================================================================
\begin{abstract}
Dengue causes recurring epidemics in Sri Lanka, and district-level forecasts three
weeks ahead give health services time to act. Spatio-temporal graph neural networks
(ST-GNNs) and SEIR-informed neural models have been proposed for this task, but both
predict the case level directly and are evaluated without the naive persistence
forecast, on a benchmark array that we show leaks future data. On this target last
week's count explains $r^2=0.89$ of the next and the best climate covariate $0.03$,
so a direct head must relearn an identity map through its encoder, and several
published ST-GNNs fail to. We propose \model{} (Persistence-Anchored Gated
Estimation), a forecasting head that plugs into any spatio-temporal encoder and
predicts a gated departure from the last observed week, at the cost of one
parameter. We rebuild a leakage-free 559-week dataset from hash-verified sources,
extend it to August 2026, and evaluate five published ST-GNNs and the LSTM of
SEIR-LSTM on one pre-registered harness. \model{} repairs the three encoders that
fail as published by 5.95--16.29 validation RMSE and 23.24--34.50 test RMSE (9/9
paired runs each), and on AAGCN it beats persistence on all 9 runs on both splits. An
ablation attributes the gain to the anchor and the gate: an SEIR simulator in the
same head adds nothing, and on 85 unseen forecast starts it ties its no-SEIR twin
($\Delta=-0.03$, $p=0.67$). On those unseen weeks no learned model beats persistence
on RMSE (34.75 against 35.43 for \model{}), although \model{} has lower MAE and
SMAPE. On this target the useful prior is persistence, not transmission dynamics.
Code and data: \codeurl.
\end{abstract}
